\chapter{Complexity — Normalize}
%%% generated by the blueprint pipeline from TCSlib.Complexity.PolyHierarchy.Normalize
\section{Overview}

Declarations of the Lean module \texttt{TCSlib.Complexity.PolyHierarchy.Normalize}.

\section{Declarations}

\begin{definition}[Iterated extension relation]
\lean{Complexity.PolyHierarchy.Extends}
\label{Complexity.PolyHierarchy.Extends}
\leanok
\uses{Turing.pairEncode}
For bit strings $w,z$ and a nonnegative integer $k$, define $z$ to extend $w$ in $k$
steps recursively: at zero steps this means $z=w$; at $k+1$ steps it means there is a
bit string $v$ such that $z$ extends the self-delimiting pairing of $w$ and $v$ in $k$
steps.
\end{definition}

\begin{theorem}[Recovery from Iterated Extensions]
\lean{Complexity.PolyHierarchy.Extends.iterate_pairFstD}
\label{Complexity.PolyHierarchy.Extends.iterate_pairFstD}
\leanok
\uses{Complexity.PolyHierarchy.Extends, Complexity.pairFstD, Complexity.pairFstD_pairEncode}
\difficulty{4}
Let $k\in\mathbb{N}$ and let $w,z$ be finite binary strings. If $z$ is obtained from $w$
by $k$ successive extensions, where each extension replaces a string $u$ by the pair
encoding of $u$ and some binary string, then applying the total first projection to $z$
exactly $k$ times returns $w$.
\end{theorem}

\begin{theorem}[Transfer of one quantified step]
\lean{Complexity.PolyHierarchy.qStep_transfer}
\label{Complexity.PolyHierarchy.qStep_transfer}
\leanok
\uses{Complexity.qStep}
\difficulty{5}
Let $b$ be a Boolean value, and let $B,m$ be nonnegative integers. Suppose $d$ maps bit
strings of length $B$ to bit strings of length $m$, every length-$m$ bit string is the
image under $d$ of some length-$B$ bit string, and for every length-$B$ string $v$,
$Q'(v)$ holds exactly when $Q(d(v))$ holds. Then the statement that there exists a
length-$B$ string satisfying $Q'$ when $b$ is true, or that every length-$B$ string
satisfies $Q'$ when $b$ is false, is equivalent to the corresponding existential or
universal statement for length-$m$ strings and predicate $Q$.
\end{theorem}

\begin{theorem}[Transfer of alternating quantifiers]
\lean{Complexity.PolyHierarchy.altQuant_transfer}
\label{Complexity.PolyHierarchy.altQuant_transfer}
\leanok
\uses{Complexity.PolyHierarchy.Extends, Complexity.PolyHierarchy.qStep_transfer, Complexity.altQuant, Complexity.altQuant_succ, Turing.pairEncode}
\difficulty{5}
Let $V$ and $V'$ be languages of bit strings, let $B,m\in\mathbb N$, and let $d$ map bit
strings to bit strings. Suppose $d(v)$ has length $m$ whenever $v$ has length $B$, and
every bit string of length $m$ equals $d(v)$ for some bit string $v$ of length $B$. Let
$(F_n)_{n\in\mathbb N}$ be a family of maps on bit strings satisfying $F_{n+1}(\langle
w,v\rangle)=\langle F_n(w),d(v)\rangle$ for all $n,w,v$, where $\langle w,v\rangle$ is
the encoding that doubles each bit of $w$, appends $01$, and then appends $v$. For every
$k\in\mathbb N$, $b\in\{\mathrm{false},\mathrm{true}\}$, $n\in\mathbb N$, and bit string
$w$, suppose every string $z$ obtained from $w$ by $k$ successive pair extensions
satisfies $z\in V'\iff F_{n+k}(z)\in V$. Then the depth-$k$ alternating-quantifier
condition on input $w$, starting with an existential quantifier if $b=\mathrm{true}$ and
a universal quantifier if $b=\mathrm{false}$, with each quantifier ranging over strings
of length $B$ and with membership in $V'$ as its depth-zero condition, is equivalent to
the corresponding condition on input $F_n(w)$ with each quantifier ranging over strings
of length $m$ and membership in $V$ at depth zero.
\end{theorem}

\begin{definition}[Decoded and padded function input]
\lean{Complexity.PolyHierarchy.nBase}
\label{Complexity.PolyHierarchy.nBase}
\leanok
\uses{Complexity.PolyHierarchy.padDecode, Complexity.pairFstD, Complexity.pairSndD, Turing.pairEncode}
Given a function $f$ from bit strings to bit strings and bit strings $s,w$, form an
input to $f$ by pairing the first component decoded from $w$ with the result of padding
the second component decoded from $w$ using the first component decoded from $s$.
\end{definition}

\begin{definition}[Padded second-component decoder]
\lean{Complexity.PolyHierarchy.nDec}
\label{Complexity.PolyHierarchy.nDec}
\leanok
\uses{Complexity.PolyHierarchy.padDecode, Complexity.pairSndD}
For bit strings $s$ and $v$, first take the second component of the pair parse of $s$,
using the empty string if parsing fails. Apply the padded decoder to that string and
$v$: it returns the prefix before the last $1$ in $v$ if that prefix has the same length
as the first string, and otherwise returns the first string; if there is no $1$ in $v$,
it returns the first string.
\end{definition}

\begin{definition}[Unary paired length encoding]
\lean{Complexity.PolyHierarchy.nSide}
\label{Complexity.PolyHierarchy.nSide}
\leanok
\uses{Turing.pairEncode}
Let $\ell_1,\ell_2 : \{0,1\}^* \to \mathbb{N}$ be functions, and let $y$ be a finite bit
string. The encoding of $y$ is the pair encoding of the unary strings of lengths
$\ell_1(y)$ and $\ell_2(y)$, respectively.
\end{definition}

\begin{definition}[Iterated tuple map]
\lean{Complexity.PolyHierarchy.nMap}
\label{Complexity.PolyHierarchy.nMap}
\leanok
\uses{Complexity.PolyHierarchy.nBase, Complexity.PolyHierarchy.nDec, Complexity.PolyHierarchy.nSide, Complexity.PolyHierarchy.tupleDecode, Complexity.pairFstD}
Let $f$ be a function on bit strings, let $\ell_1,\ell_2$ assign natural numbers to bit
strings, and let $i$ be a natural number. On input $z$, form the $(i+1)$-fold iterated
first projection of $z$, where each projection takes the first component of a pair parse
and returns the empty string if parsing fails. Use this string to construct the unary
paired length encoding from $\ell_1,\ell_2$, then apply the $(i+1)$-level tuple decoder
with base map $n\mapsto f(n)$ and the padded second-component decoder, using that
encoding and $z$ as inputs.
\end{definition}

\begin{theorem}[Polynomial-time tuple transformation]
\lean{Complexity.PolyHierarchy.polyTimeComputable_nMap}
\label{Complexity.PolyHierarchy.polyTimeComputable_nMap}
\leanok
\uses{Complexity.PolyHierarchy.UnaryPT, Complexity.PolyHierarchy.nBase, Complexity.PolyHierarchy.nDec, Complexity.PolyHierarchy.nMap, Complexity.PolyHierarchy.nSide, Complexity.PolyHierarchy.polyTimeComputable_padDecode, Complexity.PolyHierarchy.polyTimeComputable_tupleDecode, Complexity.PolyTimeComputable, Complexity.PolyTimeComputable.comp, Complexity.PolyTimeComputable.pairEncode, Complexity.pairFstD, Complexity.pairSndD, Complexity.polyTimeComputable_id, Complexity.polyTimeComputable_iterate_pairFstD, Complexity.polyTimeComputable_pairFstD, Complexity.polyTimeComputable_pairSndD}
\difficulty{4}
Let $f$ be a polynomial-time computable function on bit strings, and let $\ell_1,\ell_2$
be functions from bit strings to nonnegative integers such that the unary words of
lengths $\ell_1(y)$ and $\ell_2(y)$ are polynomial-time computable from $y$. For each
$i\in\mathbb N$, the following transformation is polynomial-time computable: given $z$,
set $s=\pi_1^{i+1}(z)$, where $\pi_1$ is the total first projection of the pair
encoding; then apply the $i$-fold tuple decoder to $z$ using the base map $(s,w)\mapsto
f(\pi_1(w),\operatorname{pad}(\pi_1(s),\pi_2(w)))$, decoder
$(s,v)\mapsto\operatorname{pad}(\pi_2(s),v)$, and side word $\langle
1^{\ell_1(s)},1^{\ell_2(s)}\rangle$.
\end{theorem}

\begin{theorem}[A normal alternating class from polynomial data]
\lean{Complexity.PolyHierarchy.mem_altClass_of_normal}
\label{Complexity.PolyHierarchy.mem_altClass_of_normal}
\leanok
\uses{Complexity.P, Complexity.PolyHierarchy.Extends.iterate_pairFstD, Complexity.PolyHierarchy.UnaryPT, Complexity.PolyHierarchy.UnaryPT.bound, Complexity.PolyHierarchy.altQuant_transfer, Complexity.PolyHierarchy.exists_padDecode_eq, Complexity.PolyHierarchy.nBase, Complexity.PolyHierarchy.nDec, Complexity.PolyHierarchy.nMap, Complexity.PolyHierarchy.nSide, Complexity.PolyHierarchy.padDecode, Complexity.PolyHierarchy.polyTimeComputable_nMap, Complexity.PolyHierarchy.qStep_transfer, Complexity.PolyHierarchy.tupleDecode, Complexity.PolyHierarchy.tupleDecode_succ_pairEncode, Complexity.PolyTimeComputable, Complexity.altClass, Complexity.altQuant, Complexity.altQuant_succ, Complexity.pairFstD, Complexity.pairFstD_pairEncode, Complexity.pairSndD, Complexity.preimage_mem_P, Complexity.qStep, Turing.pairEncode}
\difficulty{6}
Let $b\in\{\mathrm{false},\mathrm{true}\}$, let $i\in\mathbb N$, and let $V$ be a
language in polynomial time. Suppose $f$ is a polynomial-time computable function on bit
strings, and $\ell_1,\ell_2$ are functions from bit strings to natural numbers for which
the unary words $1^{\ell_1(y)}$ and $1^{\ell_2(y)}$ are polynomial-time computable from
$y$. Then the language of strings $y$ for which a quantifier of type $b$ over strings
$U$ of length $\ell_1(y)$ makes the depth-$i$ alternating-quantifier statement over $V$,
with initial type opposite to $b$, block length $\ell_2(y)$, and input $f(\langle
y,U\rangle)$ hold belongs to the alternating class at level $i+1$ with initial type $b$.
\end{theorem}

\begin{theorem}[A uniform alternating class]
\lean{Complexity.PolyHierarchy.mem_altClass_of_uniform}
\label{Complexity.PolyHierarchy.mem_altClass_of_uniform}
\leanok
\uses{Complexity.P, Complexity.PolyHierarchy.UnaryPT, Complexity.PolyHierarchy.mem_altClass_of_normal, Complexity.PolyTimeComputable, Complexity.PolyTimeComputable.comp, Complexity.PolyTimeComputable.pairEncode, Complexity.altClass, Complexity.altClass_zero, Complexity.altQuant, Complexity.altQuant_succ, Complexity.pairFstD, Complexity.pairSndD, Complexity.polyTimeComputable_pairFstD, Complexity.polyTimeComputable_pairSndD, Complexity.preimage_mem_P, Turing.pairEncode}
\difficulty{4}
Let $b\in\{\mathrm{false},\mathrm{true}\}$ and $i\in\mathbb N$. Let $V$ be a language in
polynomial time, let $f$ be a polynomial-time computable function on bit strings, and
let $\ell$ map bit strings to natural numbers with $y\mapsto 1^{\ell(y)}$
polynomial-time computable. Then the language of strings $y$ for which the depth-$i$
alternating-quantifier statement over $V$, with initial type $b$, block length
$\ell(y)$, and input $f(y)$ holds belongs to the alternating class at level $i$ with
initial type $b$.
\end{theorem}

\begin{theorem}[Polynomial-time preimages of alternating classes]
\lean{Complexity.PolyHierarchy.preimage_mem_altClass}
\label{Complexity.PolyHierarchy.preimage_mem_altClass}
\leanok
\uses{Complexity.PolyHierarchy.mem_altClass_of_uniform, Complexity.PolyHierarchy.unaryPT_poly, Complexity.PolyTimeComputable, Complexity.altClass}
\difficulty{4}
Let $b\in\{\mathrm{false},\mathrm{true}\}$, let $i\in\mathbb N$, and let $L$ be a
language in the alternating class of level $i$ whose first quantifier is specified by
$b$. If $f$ is a polynomial-time computable function on bit strings, then the preimage
$\{x:f(x)\in L\}$ also belongs to that same alternating class.
\end{theorem}
